% =============================================================================
%  论文阅读报告模板
%  编译：xelatex document.tex   （需要编译两次，第二次才生成正确的交叉引用）
%  图片：全部放在同目录的 pic/ 下，用 \includegraphics{pic/文件名}（不带扩展名）
%  结构：固定 8 节，节标题不要改名——归档脚本靠节标题定位内容，
%        详见 docs/02-阅读报告写作规范.md
% =============================================================================
\documentclass[12pt,a4paper]{article}

% ---------- 中文与版式 ----------
\usepackage[UTF8]{ctex}          % 中文排版（用 xelatex 编译）
\usepackage{geometry}
\geometry{a4paper,left=2.5cm,right=2.5cm,top=2.5cm,bottom=2.5cm}
\usepackage{setspace}
%\usepackage{parskip}             % 取消“首行缩进”，按需开启

% ---------- 数学 ----------
\usepackage{amsmath,amssymb}

% ---------- 图表 ----------
\usepackage{graphicx}
\usepackage{float}               % 允许 [H] 就地放置
\usepackage{caption}
\usepackage{subcaption}
\usepackage{booktabs}            % 三线表
\usepackage{multirow}

% ---------- 引用与链接（hyperref 尽量放最后，cleveref 放在 hyperref 之后）----------
\usepackage{url}
\usepackage{hyperref}
\usepackage{cleveref}

% ---------- 图表题注编号风格（按导师/期刊要求改这里即可）----------
\captionsetup[figure]{name=Fig.}
\renewcommand{\tablename}{Table}

% ---------- 标题（可以不写个人信息，归档时不使用）----------
%\title{论文阅读报告}
%\author{姓名}
%\date{\today}

\begin{document}
	%\maketitle

	% -------------------------------------------------------------------
	% 1. 阅读文献来源
	%    一行完整引用信息（作者、题目、期刊/会议、卷期页、年份、DOI）。
	%    第一行会被归档脚本抓取为卡片的「文献来源」。
	%    接着贴一张论文首页截图（pic/screenshot001）。
	% -------------------------------------------------------------------
	\section{阅读文献来源}
	\noindent
	% 示例：
	% A. Author, B. Author and C. Author, "Paper Title," in Journal Name, vol. 33, pp. 1-15, 2024, doi: 10.xxxx/xxxx.
	\begin{figure}[H]
		\centering
		\includegraphics[width=1\linewidth]{pic/screenshot001}
	\end{figure}

	% -------------------------------------------------------------------
	% 2. 论文待研究科学问题
	%    只写「问题」与「结论」，一段话讲清：针对什么场景/痛点、提出什么方法、
	%    关键设计点、在哪几个数据集上、达到了什么量化结果。
	%    这一段写成摘要式压缩，是为了让半年后的自己只看这段就能回忆全文。
	% -------------------------------------------------------------------
	\section{论文待研究科学问题}

	% -------------------------------------------------------------------
	% 3. 作者研究动机
	%    现有方法为什么不够用（分点写清 gap）→ 作者想解决什么 → 期望带来什么收益。
	% -------------------------------------------------------------------
	\section{作者研究动机}

	% -------------------------------------------------------------------
	% 4. 论文工作的主要内容
	%    报告的主体。建议顺序：
	%      (1) 一句话定位方法（解决什么、用什么思路）
	%      (2) 整体框架（配图 + 逐模块说明）
	%      (3) \subsection{基础知识铺垫}：读懂本文必需的前置概念，
	%          自己给自己讲一遍，不要抄论文的 related work
	%      (4) 关键公式：每个公式后必须写清符号含义与设计动机
	%    篇幅参考：占全文 50%~60%
	% -------------------------------------------------------------------
	\section{论文工作的主要内容}

	\subsection{基础知识铺垫}

	% -------------------------------------------------------------------
	% 5. 实验数据与评估指标
	%    数据集：名称、模态、规模、划分方式；指标：定义与取值范围。
	%    每个实验设置都要能回答「为什么用这个数据集/指标」。
	% -------------------------------------------------------------------
	\section{实验数据与评估指标}
	\subsection{实验数据}
	\noindent
	\subsection{评估指标}
	\noindent

	% -------------------------------------------------------------------
	% 6. 实验结果
	%    主实验（和谁比、差多少、为什么）→ 消融（哪个模块贡献大）→
	%    可视化/案例分析。数字要抄全，结论要落到「说明什么」。
	% -------------------------------------------------------------------
	\section{实验结果}

	% -------------------------------------------------------------------
	% 7. 作者工作评价（优点、缺点、后续工作要点）
	%    站在作者视角：这篇文章好在哪、作者自己承认的局限、
	%    以及作者提出的 future work。
	% -------------------------------------------------------------------
	\section{作者工作评价（优点、缺点、后续工作要点）}

	% -------------------------------------------------------------------
	% 8. 自己对该论文的评价（优点、缺点、可改进方案）
	%    整份报告最有价值的一节，也是归档后卡片里的「自评」。
	%    可改进方案要写成可执行的设计：改哪个模块 → 改成什么 →
	%    为什么能解决 → 公式（含符号与原理）→ 预期收益。
	%    汇报时这一节通常就是「结合自己课题能怎么用」的落点。
	% -------------------------------------------------------------------
	\section{自己对该论文的评价（优点、缺点、可改进方案）}

\end{document}
